We classify every rational periodic orbit of
$H_f(x,y)=(y,f(y)-x)$ for every monic cubic
$f(t)=t^3+bt^2+ct+a$ with $a,b,c\in\Z$.
The difficulty is to control the coexistence of different cycles
uniformly in all three coefficients, rather than to bound the length
of a single cycle. Rational periodic points are integral, and the
extrema of the entire periodic coordinate set give an integral third
root of a cubic secant remainder. This reduces all sufficiently large
diameters to ninety graphs over $\Z[D]$; an exact finite complement
contains 9,373 explicitly bounded parameter tuples. We give the full
graph procedure and its result tables, with no period or original-coefficient
cutoff. Every cycle uses at most three distinct coordinate
values and has least period $1$, $2$, $3$, $4$, or $6$. Seven exact
word and polynomial templates describe all cycles. There are at most
eleven rational periodic points in total, with equality precisely when
$f(t)=(t-h)^3-5(t-h)+2h$ for an integer $h$.
The coordinate set of all coexisting cycles can nevertheless contain
more than three values. The theorem uses ordinary iteration and counts
points without identifying the two orientations of a cycle.
